% ========================================
% INFORMAÇÕES DA TESE - UNESP
% ========================================

% ----- DADOS PESSOAIS -----
\newcommand{\autor}{Luiz André Colonetti Bet}
\newcommand{\autorcitacao}{BET, Luiz André Colonetti.}
\newcommand{\autorgenero}{m}

\def\curso{geografia}
\def\grau{Doutor}

% ----- ORIENTADOR -----
\newcommand{\orientador}{José Gilberto de Souza}
\newcommand{\orientadortitulo}{Prof. Dr.}
\newcommand{\orientadorgenero}{m}

% ----- COORIENTADOR (deixe vazio se não houver) -----
\newcommand{\coorientador}{Somente um teste}
\newcommand{\coorientadortitulo}{Profa. Dra.}
\newcommand{\coorientadorgenero}{f}

% ----- CAMPUS -----
\newcommand{\cidade}{Rio Claro}
\newcommand{\estado}{SP}
\newcommand{\campus}{\textit{Câmpus} de \cidade}
\newcommand{\pais}{Brasil}

\newcommand{\ano}{\the\year}
\newcommand{\mes}{Julho}

% ----- DATA DA DEFESA (usando datetime2) -----
\DTMsavedate{datadefesa}{2026-07-25}
\newcommand{\datadefesaextenso}{\DTMusedate{datadefesa}}

% ----- TÍTULO E SUBTÍTULO -----
\newcommand{\titulo}{Imperialismo e a forma-valor:}
\newcommand{\subtitulo}{China, concorrência destrutiva e a queda tendencial da taxa de lucro}
\newcommand{\tituloingles}{Imperialism and the Value-Form:}
\newcommand{\subtituloingles}{China, Destructive Competition and the Tendential Fall in the Rate of Profit}

% ----- PALAVRAS-CHAVE (PORTUGUÊS) -----
\newcommand{\palavrachaveum}{Geografia}
\newcommand{\palavrachavedois}{Imperialismo}
\newcommand{\palavrachavetres}{China}
\newcommand{\palavrachavequatro}{Taxa de lucro}
\newcommand{\palavrachavecinco}{}  % Deixe vazio se não houver

% ----- KEYWORDS (INGLÊS) -----
\newcommand{\keywordum}{Geography}
\newcommand{\keyworddois}{Imperialism}
\newcommand{\keywordtres}{China}
\newcommand{\keywordquatro}{Rate of profit}
\newcommand{\keywordcinco}{}  % Deixe vazio se não houver

% ----- BANCA EXAMINADORA -----
\newcommand{\membrodois}{Nome do Membro 2}
\newcommand{\titulomembrodois}{Prof. Dr.}
\newcommand{\instmembrodois}{Universidade do Exemplo}
\newcommand{\siglamembrodois}{EXEMPLO}

\newcommand{\membrotres}{Nome do Membro 3}
\newcommand{\titulomembrotres}{Profa. Dra.}
\newcommand{\instmembrotres}{Universidade do Exemplo}
\newcommand{\siglamembrotres}{EXEMPLO}

\newcommand{\membroquatro}{Nome do Membro 4}
\newcommand{\titulomembroquatro}{Prof. Dr.}
\newcommand{\instmembroquatro}{Universidade do Exemplo}
\newcommand{\siglamembroquatro}{EXEMPLO}

\newcommand{\membrocinco}{Nome do Membro 5}
\newcommand{\titulomembrocinco}{Profa. Dra.}
\newcommand{\instmembrocinco}{Universidade do Exemplo}
\newcommand{\siglamembrocinco}{EXEMPLO}

% ========================================
% INFORMAÇÕES INSTITUCIONAIS (NÃO ALTERAR)
% ========================================

\newcommand{\universidade}{Universidade Estadual Paulista}
\newcommand{\universidadenome}{Júlio de Mesquita Filho}
\newcommand{\universidadeabreviada}{UNESP}

% ========================================
% COMANDOS DERIVADOS (NÃO ALTERAR)
% ========================================

\ifthenelse{\equal{\curso}{geografia}}{%
    \newcommand{\institutenome}{Instituto de Geociências e Ciências Exatas}
    \newcommand{\institutoabreviado}{IGCE}
    \newcommand{\programa}{Programa de Pós-Graduação em Geografia}
    \newcommand{\area}{Geografia}
    \newcommand{\linhapesquisa}{Linha de Pesquisa em Geografia}
}{%
    \ifthenelse{\equal{\curso}{educacao}}{%
        \newcommand{\institutenome}{Instituto de Biociências}
        \newcommand{\institutoabreviado}{IB}
        \newcommand{\programa}{Programa de Pós-Graduação em Educação}
        \newcommand{\area}{Educação}
        \newcommand{\linhapesquisa}{Linha de Pesquisa em Educação}
    }{%
        \newcommand{\institutenome}{Instituto}
        \newcommand{\institutoabreviado}{-}
        \newcommand{\programa}{Programa de Pós-Graduação}
        \newcommand{\area}{Área de Concentração}
        \newcommand{\linhapesquisa}{Linha de Pesquisa}
    }%
}

\ifthenelse{\equal{\grau}{Mestre}}{%
    \newcommand{\grautxt}{Mestre}
    \newcommand{\tipotxt}{Dissertação}
    \newcommand{\graupor}{Mestrado}
}{%
    \newcommand{\grautxt}{Doutor}
    \newcommand{\tipotxt}{Tese}
    \newcommand{\graupor}{Doutorado}
}

\ifthenelse{\equal{\autorgenero}{m}}{%
    \ifthenelse{\equal{\grau}{Mestre}}{%
        \newcommand{\grauautor}{Mestre}
    }{%
        \newcommand{\grauautor}{Doutor}
    }%
}{%
    \ifthenelse{\equal{\grau}{Mestre}}{%
        \newcommand{\grauautor}{Mestra}
    }{%
        \newcommand{\grauautor}{Doutora}
    }%
}

\ifthenelse{\equal{\orientadorgenero}{m}}{%
    \newcommand{\orientadorfuncao}{Orientador}
    \newcommand{\orientadorartigo}{o}
    \newcommand{\orientadorpreposicao}{do}
}{%
    \newcommand{\orientadorfuncao}{Orientadora}
    \newcommand{\orientadorartigo}{a}
    \newcommand{\orientadorpreposicao}{da}
}

\newcommand{\orientadornome}{\orientadortitulo\ \orientador}

\ifthenelse{\equal{\coorientadorgenero}{m}}{%
    \newcommand{\coorientadorfuncao}{Coorientador}
    \newcommand{\coorientadorartigo}{o}
    \newcommand{\coorientadorpreposicao}{do}
}{%
    \newcommand{\coorientadorfuncao}{Coorientadora}
    \newcommand{\coorientadorartigo}{a}
    \newcommand{\coorientadorpreposicao}{da}
}

\ifthenelse{\equal{\coorientador}{}}{%
    \newcommand{\temcoorientador}{false}
    \newcommand{\coorientadornome}{}
}{%
    \newcommand{\temcoorientador}{true}
    \newcommand{\coorientadornome}{\coorientadortitulo\ \coorientador}
}